The emergence and evolution of morality represents one of the most enduring puzzles in evolutionary biology and social sciences~\citep{haidt2007moral, greene2013moral}. From an evolutionary standpoint, natural selection should favor individuals who maximize their reproductive success, often through selfish behaviors that increase resource acquisition at others' expense~\citep{dawkins1976selfish}. Yet, humans and some other species have evolved complex moral systems that frequently promote cooperation, altruism, and other prosocial behaviors that can seemingly contradict individual fitness maximization~\citep{tomasello2016natural}. This apparent contradiction presents a profound scientific question: Under what conditions does morality provide an evolutionary advantage?

Prior research has approached this question through evolutionary game theory~\citep{nowak2006five, axelrod1981evolution}, anthropological fieldwork on moral universals and cultural variation~\citep{henrich2015secret, curry2019good}, and biological mechanisms including kin selection~\citep{hamilton1964genetical}, reciprocal altruism~\citep{trivers1971evolution}, and group selection~\citep{wilson2007rethinking}. Moral frameworks have further identified patterns such as the expanding circle of moral concern~\citep{singer1981expanding} and fundamental moral dimensions~\citep{haidt2007moral}. While these approaches have yielded valuable insights, they share a fundamental limitation: they must abstract away the cognitive processes underlying moral behaviour, encoding agent strategies as fixed rules or payoff matrices. This prevents investigation of how cognitive factors---such as the ability to identify others' moral dispositions, memory of past interactions, and communication bandwidth---interact with environmental pressures to shape moral evolution.

Recent advances in LLMs present a novel methodological opportunity to address these limitations. LLM-based agent simulations can model entities with sophisticated cognitive architectures---including values, memory, perception, reasoning, and social dynamics---that generate emergent, complex behaviors~\citep{park2023generative, horton2023large}. This simulation paradigm allows us to observe interactions between moral cognition, behavior, and evolutionary outcomes under controlled conditions while providing rich, realistic psychological details that surpass traditional agent-based models~\citep{aher2023using}.
Nevertheless, the existing methods barely support research on morality evolution.

To address the gaps, our research advances the study of human evolution through three primary contributions:
1) Methodological: We pioneer a computational approach using LLM-based agents to bring psychological realism to evolutionary simulations, overcoming the limitations of traditional abstract models;
2) Programmatic: We release \agentframe, an extensible agent framework, and \envframe, a versatile environment platform. Together, they enable multifaceted inquiry into social evolution—from norm formation to inter-group conflict—through both long-term and scenario-based simulation;
3) Empirical: We conduct experiments across multiple settings showing that cooperation is consistently favoured by selection, with universal and reciprocal morality exhibiting the most stable survival. We discover emergent mechanisms---including a self-purging effect among selfish agents and the decisive role of the cost of moral judgment---that are difficult to study with traditional approaches. We validate robustness across multiple LLM backbones, architecture ablations, and prompt sensitivity analyses.

\begin{table*}[t]
\setlength{\tabcolsep}{4pt}
\centering
\caption{Comparison of Existing LLM-based Simulation Frameworks for Social Science.}
% \footnotesize
\small
\begin{tabularx}{\textwidth}{@{}p{3cm} c c c c c c@{}}
\toprule
\textbf{Environment} & \textbf{Morality} & \textbf{Memory} & \textbf{Team Formation} & \textbf{Competition} & \textbf{Evolution} & \textbf{Engine} \\
\midrule
\citet{park2023generative} & \xmark & Short/Long-term &  \xmark & \xmark & \xmark & LLM \\
\citet{li2023camel} & \xmark & Undefined & Multi-round Negotiation & \xmark & \xmark & LLM \\
\citet{horton2023large} & \xmark & Undefined & \xmark & Explicit & \xmark & LLM \\
\citet{aher2023using} & \xmark & Short-term & \xmark & \xmark & \xmark & LLM \\
\citet{wang2023humanoid} & \xmark & Short-term & \xmark & \xmark & \xmark & LLM \\
\citet{huang2024adasociety} & \xmark & Short-term & Multi-round Negotiation & Implicit & \cmark & LLM/RL \\
\citet{wang2024simulating} & \xmark & Short-term & \xmark & \xmark & \xmark & LLM \\
\citet{piao2025agentsociety} & \xmark &  Short/Long-term & Task-dependent & Task-dependent & \cmark & LLM \\
\citet{guan2024richelieu} & \xmark & Short/Long-term & Multi-round Negotiation & Explicit & \cmark & LLM \\
\citet{dai2024artificial} & \xmark & Short/Long-term & Negotiation-based & Explicit & \cmark & LLM \\
\midrule
Ours & \cmark & Short/Long-term & Multi-round Negotiation & Explicit & \cmark & LLM \\
\bottomrule
\end{tabularx}
\label{tab:comparison}
\end{table*}